\documentclass[11pt,a4paper]{article}
\usepackage[T1]{fontenc}
\usepackage[utf8]{inputenc}
\usepackage[french]{babel}
\usepackage{lmodern,microtype}
\usepackage[left=22mm,right=22mm,top=23mm,bottom=24mm]{geometry}
\usepackage{graphicx,booktabs,array,tabularx,longtable}
\usepackage{xcolor,tikz}
\usepackage{enumitem,titlesec}
\usepackage{hyperref}
\hypersetup{colorlinks=true,linkcolor=sfigBlue,citecolor=sfigBlue,urlcolor=sfigBlue,pdfauthor={Projet LANCE},pdflang={fr-FR}}
\urlstyle{same}
\setlength{\parindent}{0pt}
\setlength{\parskip}{5pt}
\setlength{\emergencystretch}{2em}
\setlist{nosep,leftmargin=5mm,topsep=3pt}
\makeatletter
\long\def\@makecaption#1#2{\vskip 7pt\begingroup\small\textbf{#1} -- #2\par\endgroup\vskip 3pt}
\def\ps@lance{%
  \def\@oddhead{\small\sffamily LANCE / Recherche\hfil 29 septembre 2026}%
  \let\@evenhead\@oddhead
  \def\@oddfoot{\small\color{sfigMuted}Diagnostic logiciel et protocole proposé\hfil\thepage}%
  \let\@evenfoot\@oddfoot
}
\makeatother
\def\UrlBreaks{\do\/\do\-\do\_\do\.\do\?\do\&\do\=\do\:\do\,}

\newcolumntype{Y}{>{\raggedright\arraybackslash}X}
\newcommand{\code}[1]{{\small\ttfamily\path{#1}}}
\newcommand{\evidence}[1]{\par\smallskip{\small\color{sfigMuted}\textbf{Preuve / limite.} #1}\par\smallskip}
\newcommand{\status}[1]{{\small\sffamily\color{sfigRaspberry}#1}}
\titleformat{\section}{\Large\sffamily\bfseries\color{sfigInk}}{\thesection}{0.7em}{}
\titleformat{\subsection}{\large\sffamily\bfseries\color{sfigInk}}{\thesubsection}{0.7em}{}
\titlespacing*{\section}{0pt}{16pt}{7pt}
\titlespacing*{\subsection}{0pt}{11pt}{5pt}
\pagestyle{lance}
\setlength{\headheight}{14pt}

% Reusable TikZ styles for sober scientific figures.
% Required packages: xcolor, tikz.
% Required TikZ libraries: arrows.meta, positioning, fit, calc.

\usetikzlibrary{arrows.meta,positioning,fit,calc}

\definecolor{sfigInk}{HTML}{1F242C}
\definecolor{sfigMuted}{HTML}{5C626C}
\definecolor{sfigLine}{HTML}{D9DEE5}
\definecolor{sfigNeutral}{HTML}{F7F8FA}
\definecolor{sfigBlue}{HTML}{1D5C96}
\definecolor{sfigBlueLine}{HTML}{93B1CD}
\definecolor{sfigBlueFill}{HTML}{EFF6FC}
\definecolor{sfigRaspberry}{HTML}{AF0032}
\definecolor{sfigRoseLine}{HTML}{DB8CA3}
\definecolor{sfigRoseFill}{HTML}{FDF0F4}

\tikzset{
  sfig diagram/.style={
    >=Latex,
    node distance=10mm and 12mm,
    every node/.style={font=\sffamily\small,text=sfigInk}
  },
  sfig base/.style={
    draw=sfigLine,
    fill=white,
    rounded corners=2pt,
    line width=0.55pt,
    inner xsep=6pt,
    inner ysep=4pt,
    align=center
  },
  sfig input/.style={
    sfig base,
    draw=sfigBlueLine,
    fill=sfigBlueFill
  },
  sfig process/.style={
    sfig base,
    draw=sfigRoseLine,
    fill=sfigRoseFill
  },
  sfig neutral/.style={
    sfig base,
    fill=sfigNeutral
  },
  sfig focal/.style={
    sfig base,
    draw=sfigRaspberry,
    fill=sfigRoseFill,
    line width=0.75pt,
    font=\sffamily\small\bfseries,
    text=sfigRaspberry
  },
  sfig output/.style={
    sfig neutral,
    font=\sffamily\small\bfseries
  },
  sfig group/.style={
    draw=sfigLine,
    fill=none,
    dashed,
    rounded corners=3pt,
    line width=0.5pt,
    inner sep=6pt
  },
  sfig flow/.style={
    -{Latex[length=2mm,width=1.2mm]},
    draw=sfigMuted,
    line width=0.6pt
  },
  sfig blue flow/.style={
    -{Latex[length=2mm,width=1.2mm]},
    draw=sfigBlue,
    line width=0.68pt
  },
  sfig primary flow/.style={
    -{Latex[length=2mm,width=1.2mm]},
    draw=sfigRaspberry,
    line width=0.78pt
  },
  sfig blocked/.style={
    draw=sfigRaspberry,
    dashed,
    line width=0.68pt
  },
  sfig annotation/.style={
    font=\sffamily\scriptsize,
    text=sfigMuted,
    align=center
  },
  sfig formula/.style={
    font=\small,
    text=sfigInk,
    align=center
  }
}

\hypersetup{pdftitle={LANCE : couverture, contexte et apport mesurable des LLM}}
\begin{document}
{\sffamily\small\color{sfigMuted}ÉTUDE DE RECHERCHE / RAPPORT PRINCIPAL}\par
{\sffamily\LARGE\bfseries LANCE : mesurer l'apport des LLM}\par
{\sffamily\large Couverture du benchmark, contexte du pipeline\par et protocole de comparaison}\par
\medskip
\status{Exploration et protocole proposés / 29 septembre 2026}\par
\smallskip
\textbf{Question.} Dans quelles situations un modèle de langage améliore-t-il
un audit IoT autorisé, à outils et budgets comparables à une automatisation
conditionnelle bien conçue ?\par
\textbf{État de la recherche.} Trois inspections GPT et une contrelecture
croisée ont identifié des limites concrètes. Huit probes synthétiques ont été
rejouées localement ; 381 tests logiciels ciblés passent. Aucun résultat de
campagne ne démontre encore un gain technique ou humain du modèle.

\section{Lecture d'ensemble}
Le projet dispose déjà d'un socle de collecte, de contrats de preuve et d'une
comparaison D1/A1 bornée. La prochaine étape consiste à rendre les mesures et
l'attribution suffisamment fiables pour qu'un résultat positif, négatif ou
indéterminé ait un sens. Le tableau~\ref{tab:axes} résume les trois études.

\begin{table}[htbp]\centering\small
\begin{tabularx}{\linewidth}{@{}p{25mm}YY@{}}\toprule
Axe & Faits établis dans le périmètre inspecté & Question encore ouverte \\\midrule
Benchmark & 29 scénarios ; 288 vulnérabilités attendues ; aucun scénario officiel entièrement sain. Défauts locaux de taux de tentative et d'admission de preuves. & Quelles propriétés sont réellement déployées, observées et prouvées ? Quelle fréquence d'erreur en campagne ? \\
Contexte & Certaines sorties coupées avant archivage ; historique outils renvoyé sans budget d'entrée ; projections parfois sans récupération. & Quel effet sur les décisions, le rappel, la fiabilité et les ressources ? \\
Apport LLM & Scanner partagé riche en connaissances des simulateurs ; comparaison disponible sur les phases 3--4, inventaire fourni, profil full. & Quel gain marginal du modèle et quel mécanisme le produit ? \\\bottomrule
\end{tabularx}
\caption{Faits de logiciel et de corpus ; aucune colonne ne rapporte une performance de campagne.}\label{tab:axes}
\end{table}

Le présent rapport fournit la synthèse et le protocole de décision. Deux
rapports spécialisés détaillent les preuves :
\href{../benchmark-coverage/report.pdf}{\emph{Couverture et validité du benchmark}}
et \href{../context-scalability/report.pdf}{\emph{Gestion du contexte du pipeline}}.
Les notes d'inspection et les références de code restent dans les mêmes dossiers.

\evidence{Révision inspectée : \code{45c8dab892ecb4e5f513813387090e0cded1b85d},
avec une réorganisation documentaire locale préexistante. Les changements de
cette livraison portent sur les documents, les figures et les reproductions.
Les corrections du logiciel décrites ci-dessous restent proposées.}

\section{Ce que compare aujourd'hui D1/A1}
\subsection{Un ajout au socle déterministe}
Les deux politiques utilisent un inventaire public commun et le même scanner
initial. D1 produit les findings par règles, puis suit des plans de vérification
avec une reprise transitoire et un complément HTTP bornés. A1 ajoute des
micro-agents en phase 3 et des décisions modèle en phase 4. La figure~\ref{fig:paired}
représente ce périmètre implémenté ; elle ne décrit pas encore un audit complet
avec reconnaissance, intrusion et rédaction comparées.

\begin{figure}[htbp]\centering
\begin{tikzpicture}[sfig diagram,x=0.94cm]
\node[sfig input,text width=21mm] (inventory) at (0,0) {Inventaire\\public commun};
\node[sfig neutral,text width=24mm] (scanner) at (3.2,0) {Scanner\\partagé};
\node[sfig base,text width=26mm] (d1) at (7,1) {D1\\règles et reprises};
\node[sfig focal,text width=26mm] (a1) at (7,-1) {A1\\décisions LLM};
\node[sfig base,text width=24mm] (proof) at (10.8,0) {Contrat de\\preuve commun};
\node[sfig output,text width=22mm] (score) at (14,0) {Scores\\et ressources};
\draw[sfig flow] (inventory) -- (scanner);
\draw[sfig flow] (scanner.east) -- +(5mm,0) |- (d1.west);
\draw[sfig primary flow] (scanner.east) -- +(5mm,0) |- (a1.west);
\draw[sfig flow] (d1.east) -- +(5mm,0) |- (proof.west);
\draw[sfig primary flow] (a1.east) -- +(5mm,0) |- (proof.west);
\draw[sfig flow] (proof) -- (score);
\end{tikzpicture}
\caption{Comparaison actuellement disponible : effet marginal de la politique LLM au-dessus du scanner partagé, dans les phases 3--4. Le validateur commun assure une symétrie interne ; sa sémantique doit aussi être calibrée indépendamment.}\label{fig:paired}
\end{figure}

Un succès absolu d'A1 ne prouve pas que le modèle a causé ce succès. Le scanner
contient déjà des endpoints, identités et séquences propres aux simulateurs.
La provenance d'une action est utile pour suivre une trajectoire, mais une
ablation contrôlée est nécessaire pour attribuer un gain au modèle.

\evidence{Chemins inspectés : \code{src/agent/audit_experiment.py},
\code{src/agent/phases/analysis/run.py}, \code{src/agent/scanner.py},
\code{src/agent/phases/verification/rules.py}. La branche de mémo phase 3 compact
local décrite dans l'étude contexte est distincte : D1/A1 impose le profil full.}

\subsection{Une référence classique à renforcer sur développement}
D1 est déjà conditionnel, mais borné. Une victoire d'A1 sur une lacune connue et
réparable de D1 n'établirait pas la nécessité d'un LLM. Avant le gel, revoir
l'extraction des paramètres observables, les branchements, l'état partagé,
les identités découvertes, les contrôles négatifs et les critères d'arrêt.
Les workflows conditionnels de Nuclei illustrent des possibilités classiques
sans modèle~\cite{nuclei} ; ils ne garantissent pas une couverture IoT universelle.

Établir une matrice d'origine des connaissances pour chaque famille :
information publique fournie, observation obtenue, règle scanner, prompt,
connaissance du modèle ou oracle privé. Les deux politiques ont accès aux mêmes
informations opérationnelles autorisées. Le test réservé ne fournit aucune
réponse via les noms, rôles ou descriptions.

\section{Fiabiliser les dimensions de couverture}
\subsection{Corpus, tentative et preuve ont des dénominateurs différents}
Le corpus officiel comprend 19 scénarios de développement et 10 de test,
respectivement 252 et 36 GT positives. Les 288 GT, 249 instances de services et
67 chemins déclarés ne constituent pas des mécanismes indépendants. Les
scénarios S1--S13 regroupent 229 GT ; le test est surtout HTTP/SSH. Les variantes
historiques S1h/S4h ne s'ajoutent pas aux 29 scénarios officiels.

Une tentative sur tous les candidats retenus ne signifie pas que tous les actifs,
services ou failles ont été explorés. Conserver des populations distinctes :
actifs attendus lorsque connus, services observés, GT proposées, candidats
retenus, candidats effectivement tentés, preuves admises et GT confirmées.
Le score d'intrusion reste séparé du F1 d'audit.

\subsection{Contre-exemples locaux et portée exacte}
\begin{table}[htbp]\centering\small
\begin{tabularx}{\linewidth}{@{}p{30mm}YY@{}}\toprule
Contre-exemple & Observation reproduite & Portée du constat \\\midrule
Tentatives ambiguës & Un candidat, deux TIMEOUT marqués non tentés : taux de tentative 1,0. & Défaut du diagnostic d'entonnoir ; fréquence en campagne inconnue. \\
Contrôle S18 & Une exposition avec jeton surprivilégié partage IP/type avec un contrôle exigeant le refus d'un autre jeton limité. & Contrôle déclaré violé avec le mauvais jeton ; aucune pénalité additionnelle directe démontrée sur le F1 principal. \\
Oracle OTA & Une trace déclarant une signature valide est admise comme \code{EXPLOITED}, niveau 2. & Admission d'une observation insuffisante ; aucune signature cryptographique ni état serveur vérifié. \\
Oracle exécution & Une réponse d'aide contenant \code{uid=1000(user)} est admise au niveau 3. & Un exemple textuel est pris pour preuve ; le crédit dans le score complet n'est pas reproduit. \\\bottomrule
\end{tabularx}
\caption{Défauts et limites identifiés par inspection ou probes synthétiques, sans extrapolation à un taux d'erreur réel.}\label{tab:counterexamples}
\end{table}

Priorité proposée : séparer attribution, tentative et verdict ; définir les
contrôles par cible, service, identité, action et observation attendue ; calibrer
les oracles sur preuves positives et contre-exemples annotés indépendamment.
Un nonce réfléchi dans une réponse ne prouve pas une exécution. L'oracle doit
distinguer réflexion, exemple statique et effet attendu, avec contrôle
différentiel ou attestation de cible.

Les propriétés du laboratoire doivent être vérifiées avant l'essai. Un
précontrôle perturbateur peut modifier l'état : réinitialiser ensuite, puis
réaliser un contrôle final non perturbateur. Le hash d'un constat assure sa
traçabilité, pas la validité de son contenu. Pour les pivots, la provenance
causale d'une transition réseau reste absente : plusieurs connexions directes
ne démontrent pas un chemin multi-hop.

\section{Contexte : préserver les informations avant d'optimiser}
Le rapport de phase 6 utilise déjà des fiches indépendantes et un assemblage
déterministe. Le reprendre comme un gros contexte global ignorerait cette
amélioration. Les problèmes inspectés concernent plusieurs frontières :
\begin{itemize}
\item \textbf{Capture.} Certains outils YAML filtrent et tronquent stdout avant
que le journal ne reçoive le résultat. Une référence au journal ne récupère pas
une donnée qui n'y a jamais été conservée.
\item \textbf{Conversation.} La boucle outils renvoie l'historique intégral,
sans admission explicite selon un budget d'entrée et une réserve de sortie.
Les plafonds de tours ne garantissent pas une requête compatible avec la fenêtre.
\item \textbf{Projection.} Le contexte compact peut perdre une preuve au milieu
ou des services entiers ; certains chemins exposent uniquement la sauvegarde et
ne peuvent pas relire les références affichées.
\item \textbf{Transmission.} Des coupes de preuves et des sélections d'alternatives
par IP retirent des éléments sans signaler systématiquement leur omission.
\end{itemize}

Les cinq probes synthétiques établissent des transformations, pas leur effet
sur un score : preuve centrale perdue avec zéro entrée omise, 20 services réduits
à 6, observations tronquées, blocs de tailles $[2,2,2,14]$ et mémos incomplets
non détectés. Une sortie interrompue par \code{finish_reason=length} ne prouve
pas à elle seule une saturation de l'entrée.

La proposition est de conserver un brut attribué et versionné, puis construire
des vues déterministes portant leurs omissions et références. Fournir des
lectures ciblées, bornées et sans nouvelle sonde réseau aux politiques qui en
ont besoin. Comparer d'abord cette solution à une meilleure projection
simple : le retrieval peut augmenter coût et délai ou manquer la bonne preuve.
Les recommandations sur le chargement à la demande motivent ce test~\cite{context},
et la sensibilité à la position documentée dans d'autres tâches motive des
permutations~\cite{lost}; elles ne valident aucun modèle actuel de LANCE.

\section{Expériences proposées : isoler les contributions}
\subsection{Politiques et mécanismes}
\begin{table}[htbp]\centering\small
\begin{tabularx}{\linewidth}{@{}p{15mm}YY@{}}\toprule
Politique & Décisions & Rôle expérimental \\\midrule
D0 & Procédure fixe. & Référence simple secondaire. \\
D1 & Règles, état partagé, branchements et reprises bornées. & Référence principale sans LLM ; premier incrément disponible. \\
A0 & Plan initial LLM puis figé, branches prévues autorisées. & Isoler l'effet de nouvelles décisions après retour ; à implémenter. \\
A1 & Politique LLM révisée selon les observations. & Traitement étudié ; incrément borné disponible. \\
D1-R & Même audit D1, rédaction LLM après clôture. & Étudier restitution et effort humain ; proposé. \\\bottomrule
\end{tabularx}
\caption{Les politiques proposées dépassent la comparaison logicielle actuellement disponible ; leur présence dans ce tableau n'atteste pas leur implémentation.}\label{tab:policies}
\end{table}

\begin{table}[htbp]\centering\small
\begin{tabularx}{\linewidth}{@{}p{28mm}YY@{}}\toprule
Question & Contraste & Conclusion permise \\\midrule
Interprétation & Règles / LLM sur dossiers identiques, collecte désactivée. & Qualité d'interprétation et fidélité des preuves. \\
Audit borné & D1 / A1, mêmes inventaire, outils, état initial et plafonds. & Gain marginal au-dessus du socle partagé. \\
Adaptation & A0 / A1 après un obstacle externe reproductible. & Utilité de nouvelles décisions après retour. \\
Gestion du contexte & Même politique/modèle, stratégie de représentation seule modifiée. & Effet de la représentation et de ses coûts. \\
Restitution & D1 / D1-R, mêmes faits, étude aveugle d'auditeurs. & Temps de compréhension/correction et qualité des décisions humaines. \\\bottomrule
\end{tabularx}
\caption{Contrastes distincts pour éviter de créditer au modèle les effets de meilleurs outils, d'un index ou d'un budget différent.}\label{tab:contrasts}
\end{table}

D1-R doit conserver exactement les scores techniques de D1. Un changement
signalerait que la mesure technique dépend de la prose. Les temps humains
nécessitent une étude dédiée ; ils ne se déduisent pas des tokens ou de
l'absence d'intervention. La comparaison de plusieurs modèles répond à une
question de choix de modèle, sans remplacer D1/A1.

\subsection{Variations et ablations}
Séparer variation du format brut et recomposition des dépendances. Fournir aussi
une strate où les mêmes informations sont normalisées pour les deux politiques.
Un avantage qui disparaît après normalisation peut relever du parseur ; pour
étudier l'adaptation, introduire des informations ou obstacles qui nécessitent
une révision d'action. Les perturbations sont externes, reproductibles et
horodatées ; une panne accidentelle n'est pas requalifiée après coup.

Pour le contexte, comparer une projection actuelle, une meilleure projection
déterministe et des vues avec lectures ciblées. Garder identiques les capacités
de lecture pour isoler la projection ; retirer les lectures dans une ablation
distincte. Tester position début/milieu/fin, permutation des services, bruit,
contradictions tardives et éléments répartis entre plusieurs services. Mesurer
fidélité, attribution, abstention, omissions, lectures, temps et coût.

Compact/full change outils, budgets, obligations et rôle du modèle : ce
contraste ne constitue pas une ablation propre du contexte. Une trace figée
hors ligne ne reproduit pas les effets d'actions divergentes sur le réseau.

\section{Protocole statistique et ressources}
\subsection{Unité d'analyse}
Garder les deux bras appariés sur la même instance et le même état initial.
Réinitialiser avant chaque bras et contrebalancer l'ordre D1/A1 par blocs.
Pour la différence principale, une hiérarchie proposée est : moyenne des
répétitions dans une variante, moyenne des variantes dans une famille, puis
poids égal des familles. Avec $f$ une famille, $v$ une variante et $r$ une
répétition, poser, pour le F1 final des confirmations sur les scénarios positifs
au budget principal fixé :
\[
 d_{fvr}=F1(A1_{fvr})-F1(D1_{fvr}),\qquad
 \Delta=\frac{1}{F}\sum_f\frac{1}{V_f}\sum_v\frac{1}{R_{fv}}\sum_r d_{fvr}.
\]
Cette analyse de recherche ne modifie pas les agrégats officiels par scénario
et par split. Elle suppose une règle explicite de disponibilité des paires.
Les scénarios entièrement sains ont un F1 indéfini ; analyser séparément leurs
fausses alertes et violations, sans les inclure dans cette moyenne.
Les familles reflètent l'origine technique et les mécanismes partagés ; un
label de manifeste n'en prouve pas l'indépendance.

Estimer l'incertitude à 95\% par rééchantillonnage apparié des familles en gardant
leur structure interne, après validation du schéma de données. Sur peu de
familles, rester descriptif : le bootstrap ne crée pas de diversité. Les
principes généraux des tests appariés~\cite{dror} ne donnent pas à eux seuls la
puissance du pilote. Dimensionner le test à partir du développement, de la
variance et du gain minimal utile. Déclarer un contraste principal et les
analyses secondaires avant la lecture du test.

\subsection{Échecs, coûts et reproductibilité}
La moyenne des paires admissibles est conditionnelle aux survivants. Conserver
les essais prévus, échecs, arrêts, budgets dépassés et dépenses connues. Les
scores techniques indisponibles restent indisponibles. Ajouter séparément une
réussite opérationnelle définie avant test : résultat évaluable et niveau de
qualité atteint avant plafond, sur environnement initialement valide.

Fixer des plafonds comparables de durée murale et charge réseau : requêtes,
cibles, plages et débits. Un appel Nmap ne représente pas une seule sonde, et un
plafond de tokens n'est pas un budget commun avec D1. Séparer coût LLM marginal,
calcul/infrastructure, préparation, développement, maintenance et temps humain.
Un ratio coût/VP est indéfini si VP vaut zéro ; les dépenses restent visibles.
L'étude doit présenter les compromis qualité/coût, plutôt qu'un seul score
favorable choisi après observation~\cite{matter}.

Figer sources et diff local, règles/prompts, snapshot CVE, SDK, outils/OS,
inventaire, état initial, perturbations, endpoint sans secret ou empreinte,
version résolue du modèle et paramètres effectifs connus ou inconnus. La
configuration actuelle ne capture pas entièrement endpoint et defaults
fournisseur ; l'égalité d'un nom et du hash source ne suffit pas.

\section{Ordre de travail et décisions}
La figure~\ref{fig:workplan} indique l'ordre proposé. Le pilote peut apprendre
sur les phases 3--4 correctement vérifiables sans attendre une refonte complète
ou la preuve de tous les pivots. Ses conclusions restent descriptives.

\begin{figure}[htbp]\centering
\begin{tikzpicture}[sfig diagram,x=0.94cm]
\node[sfig focal,text width=24mm] (measure) at (0,0) {Mesures\\et oracles};
\node[sfig base,text width=24mm] (context) at (3.5,0) {Contexte\\observable};
\node[sfig base,text width=24mm] (offline) at (7,0) {Contrastes\\hors ligne};
\node[sfig base,text width=24mm] (pilot) at (10.5,0) {Pilote D1/A1\\24 essais prévus};
\node[sfig output,text width=24mm] (test) at (14,0) {Test réservé\\et mécanismes};
\draw[sfig primary flow] (measure) -- (context);
\draw[sfig flow] (context) -- (offline);
\draw[sfig flow] (offline) -- (pilot);
\draw[sfig flow] (pilot) -- (test);
\end{tikzpicture}
\caption{Séquence proposée. Les travaux logiciels, le pilote et le test confirmatoire sont des étapes futures ; aucune performance n'est dessinée. Les seuils et la taille finale sont gelés après développement, avant test.}\label{fig:workplan}
\end{figure}

\begin{samepage}
Les décisions confrontent l'estimation et son incertitude aux marges fixées
avant le test ; une moyenne favorable ne suffit pas.
\par\end{samepage}

\begin{table}[htbp]\centering\small
\begin{tabularx}{\linewidth}{@{}YY@{}}\toprule
Résultat selon critères préfixés & Décision proposée \\\midrule
Gain supérieur au minimum utile, FP/fiabilité/coût acceptables. & Conserver le LLM sur le périmètre établi ; rechercher ensuite le mécanisme. \\
D1 non inférieur selon marge fixée, moins coûteux et plus fiable. & Utiliser D1 pour les décisions ; évaluer D1-R séparément si nécessaire. \\
Gain seulement sur les cas ambigus. & Définir un routage règles vers LLM sur développement, puis évaluer le routeur et ses faux refus sur un nouveau test. \\
Interprétation meilleure hors ligne, pas de gain terrain. & Examiner collecte, contexte, exécution et preuve ; corriger sur développement puis retester. \\
FP ou fiabilité dégradés au-delà des marges. & Restreindre ou désactiver l'adaptation sur ce périmètre. \\
Incertitude trop grande ou peu de familles. & Résultat indéterminé ; ajouter diversité ou réduire la portée. \\\bottomrule
\end{tabularx}
\caption{Accepter une conclusion négative fait partie du protocole. Aucune marge numérique n'a encore été choisie.}\label{tab:decisions}
\end{table}

Il reste à fixer avant le test : gain minimal utile, marges de faux positifs et
fiabilité, budgets, règles de données manquantes, taille fixe, exclusions et
analyses secondaires. Une absence de significativité ne prouve pas
l'équivalence. Si l'utilité se situe dans la construction de règles, une autre
expérience peut comparer humain seul et humain assisté pour la maintenance,
puis exécuter le workflow déterministe ; ce n'est pas un gain d'agent à chaque audit.

\section{Validation disponible et limites de la revue}
\begin{table}[htbp]\centering\small
\begin{tabularx}{\linewidth}{@{}p{36mm}Y@{}}\toprule
Vérification & Résultat disponible au 29 septembre 2026 \\\midrule
Inventaire statique & 29 SHA de vérités terrain concordants ; totaux du corpus vérifiés avec fusion des sidecars. \\
Composition GT & 31 fichiers examinés, dont S1h/S4h ; 13 différences descriptives historiques acceptées et 18 correspondances strictes. Pas de validation \code{strict-all}. \\
Probes synthétiques & Trois probes évaluation et cinq contexte rejouées ; défauts caractérisés, fréquence terrain inconnue. \\
Tests ciblés & 381 cas réussis en 8,35 s dans un venv Python 3.12.2 / pytest 9.1.1 ; fournisseurs et outils simulés. \\
Manifeste pilote exemple & 12 paires prévues, zéro admissible, 24 bras non exécutés ; huit champs configuration non fixés et coûts inconnus. Cela décrit l'exemple uniquement. \\\bottomrule
\end{tabularx}
\caption{Validation logicielle et de définitions, distincte d'une campagne expérimentale.}\label{tab:validation}
\end{table}

Les reproductions sont conservées dans
\code{benchmarks/experiments/research-review/} : \code{inventory.py},
\code{probe-evaluation.py} et \code{probe-context.py}. Les versions, commandes,
incidents d'environnement résolus et observations sont détaillés dans
\code{research/agent-vs-automation/validation.md}. Les assertions du probe
contexte caractérisent les défauts de cette révision ; elles ne sont pas des
tests du contrat souhaité après correction.

Les métadonnées historiques examinées ne donnent pas suffisamment de
provenance pour en déduire une performance actuelle. Les données test S20--S29
sont publiques et leur indépendance historique reste non vérifiée. Préparer un
jeu réservé avec séparation réelle des fichiers et oracles ; le chemin D1/A1
sur inventaire ne bénéficie pas automatiquement de l'isolation scellée.

\begin{samepage}
La revue bibliographique est ciblée et non systématique. Les évaluations CTF
et cyber publiées motivent des précautions~\cite{autopen,cybench,cvebench}, mais
leurs scores ne sont pas comparables au F1 IoT de LANCE. Aucune performance
externe n'est reproduite ici. La contrelecture GPT a précisé l'attribution et les critères ; la contre-revue Muse préparée reste non exécutée et aucun avis
ne lui est attribué. Les questions et arbitrages figurent dans
\code{adversarial-review.md}.
\par\end{samepage}

\begin{thebibliography}{9}\small
\bibitem{matter} S. Kapoor et al. \emph{AI Agents That Matter}. arXiv:2407.01502v1, 1 juillet 2024. Revue initiale : sections 2, 5, 6 et passages d'annexe lus ; tâches programmation/navigation, pas audit IoT. Consulté le 29 septembre 2026. \url{https://arxiv.org/html/2407.01502v1}.
\bibitem{nuclei} ProjectDiscovery. \emph{Template Workflows Overview}. Documentation vivante, version/date de publication non précisées. Workflows conditionnels et contexte partagé lus ; consultation le 29 septembre 2026. \url{https://docs.projectdiscovery.io/templates/workflows/overview}.
\bibitem{context} Anthropic Engineering. \emph{Effective context engineering for AI agents}, 29 septembre 2025. Corps principal, récupération, compaction et notes structurées lus ; recommandation fournisseur, pas expérimentation LANCE. Consulté le 29 septembre 2026. \url{https://www.anthropic.com/engineering/effective-context-engineering-for-ai-agents}.
\bibitem{lost} N. F. Liu et al. \emph{Lost in the Middle: How Language Models Use Long Contexts}. arXiv:2307.03172v3, 20 novembre 2023. Résumé et historique lus ; pas article intégral. Modèles/tâches de 2023. Consulté le 29 septembre 2026. \url{https://arxiv.org/abs/2307.03172v3}.
\bibitem{dror} R. Dror et al. \emph{The Hitchhiker's Guide to Testing Statistical Significance in Natural Language Processing}. ACL 2018. Notice et sections 3.2--3.3 lues ; principes généraux, pas validation de nos grappes. Consulté le 29 septembre 2026. \url{https://aclanthology.org/P18-1128/}.
\bibitem{autopen} L. Gioacchini et al. \emph{AutoPenBench: A Vulnerability Testing Benchmark for Generative Agents}. EMNLP Industry, novembre 2025. Sections 2--5 et limitations lues ; assistance humaine et conditions à distinguer. Consulté le 29 septembre 2026. \url{https://aclanthology.org/2025.emnlp-industry.114.pdf}.
\bibitem{cybench} A. K. Zhang et al. \emph{Cybench}. ICLR 2025. Notice officielle/résumé lus dans la revue comparative ; présentation des auteurs lue dans la revue benchmark. CTF, pas mesure des contrôles sains IoT. Consulté le 29 septembre 2026. \url{https://cybench.github.io/}.
\bibitem{cvebench} Y. Zhu et al. \emph{CVE-Bench}. ICML/PMLR 267, juillet 2025. Notice et résumé lus ; exploitation web critique, pas corpus IoT exhaustif. Consulté le 29 septembre 2026. \url{https://proceedings.mlr.press/v267/zhu25i.html}.
\end{thebibliography}
\end{document}
